\documentclass[10pt,a4paper]{amsart}
\usepackage[margin=19mm]{geometry}
\usepackage{amsmath,amssymb,mathtools}
\usepackage[scr=rsfs,cal=euler]{mathalfa}
\usepackage{xfrac}
\usepackage[unicode,hidelinks,bookmarks=false]{hyperref}
\newcommand{\ie}{i.e.\ }
\newcommand{\cf}{cf.\ }
\newcommand{\resp}{respectively\ }
\newcommand{\Subsection}{\subsection}
\providecommand{\nobreakdash}{\nobreak\hbox{-}}
\setlength{\emergencystretch}{2em}
\multlinegap0pt
\usepackage{amssymb}
\usepackage{xfrac}
\usepackage{mathtools}
\usepackage[shortlabels]{enumitem}
\usepackage{float}
\newtheorem{proposition}{Proposition}
\title{There are many 5-holes}
\author{Omar Astudillo-Marb\'an, Oriol Sol\'e-Pi}
\date{Source: arXiv:2603.18484v1 (CC BY 4.0)}
\begin{document}
\maketitle

\noindent\emph{Source and licence.} There are many 5-holes. Version arXiv:2603.18484v1 (CC BY 4.0), distributed by arXiv under the Creative Commons Attribution 4.0 International licence, http://creativecommons.org/licenses/by/4.0/, as recorded in the arXiv licence metadata for this exact version and shown on the abstract page https://arxiv.org/abs/2603.18484v1. Full licence notice: \texttt{licenses/arxiv-2603.18484-CC-BY-4.0.txt}.\par\smallskip

\noindent\textit{Editorial scope.} Counted unit: the article's proposition prop:second\_structure (source0.tex lines 341--348). The article's complete original proof of it is delivered with it (source0.tex lines 350--370, one \texttt{proof} environment of 21 lines). No mathematics is written, repaired or re-derived: the statement and the proof are the article's own lines, in the article's own notation, and no retained line is substituted or re-flowed.  Retained uncounted as the source-local context the proof needs: lines 305--314; lines 371--376.  One declared adaptation: exactly 2 ancillary strings inside the retained spans are omitted (image-inclusion commands and, where the source repeats a label, the repeated label definitions) and no other line differs from the source; the figure environments, with their placement, captions and labels, are kept, and the package records the omitted strings in fidelity receipts bound to the affected bodies.  No retained line contains a citation, so the wrapper carries no bibliography and no bibliography entry.  The retained lines contain the article's own diagram code, which the wrapper typesets with the standard tikz packages; no image file is delivered and the package declares no asset dependency.  Editorial formatting only, in addition: an AMS \texttt{amsart} preamble that restates the article's own declarations and macros used by the retained lines, namely the environments \texttt{proposition} on separate counters, together with the standard packages those lines need.  The retained environments print in this document as Proposition 1 in order of appearance, whereas the article prints the same items with the numbering of its own full document.  The complete native source of the article as distributed in the arXiv e-print for this exact version is bundled unchanged as \texttt{sources/196698/source0.tex}.
\begin{proposition}\label{prop:first_structure}
For every positive integer $k$ there exists some other positive integer $K_1=K_1(k)$ such that, among any $K_1$ points to which $H$ has been assigned, we can find a subset $S$ of $k$ points for which the following properties are satisfied:
\begin{itemize}
    \item all elements of $S$ are assigned to $H$ via the same anchors $a$ and $b$ and, for any two distinct vertices $a'$ and $b'$ of $H$, all elements of $S$ lie on the same side of $\overleftrightarrow{a'b'}$;

    \item the elements of $S$ are in convex position;

    \item there is a labeling $s_1,\dots,s_{k}$ of the elements of $S$ so that $s_1\dots s_k$ is a convex polygon (in that order) and, for $1\leq i<j\leq k$, every vertex $h$ of $H$ satisfies $\measuredangle s_is_jh<\pi$.
\end{itemize}
\end{proposition}

\begin{figure}[h]
    \centering
    
    \caption{a) The quadrilateral $sp_{i-1}p_ip_{i+1}$ is a $4$-hole, and $sp_{i-1}p_ip_{i+1}s'$ is a convex pentagon. Thus, if $s''$ is the point in $\triangle sp_{i+1}s'$ with smallest non-zero Euclidean distance to $\overleftrightarrow{sp_{i+1}}$, then $sp_{i-1}p_ip_{i+1}s''$ is a $5$-hole. b) The point $q$ lies in $\angle p_{j}sp_{j+3}$. The point $q'$ has the smallest non-zero Euclidean distance to $\overleftrightarrow{p_jp_{j+3}}$ among the points in $\triangle p_jqp_{j+3}$. The pentagon $p_jp_{j+1}p_{j+2}p_{j+3}q'$ (shown in red) is a $5$-hole which induces a trustworthy assignment.}
    \label{fig:pushing_from_top}
\end{figure}

\begin{proposition}\label{prop:second_structure}
For every positive integer $k$ there exists some positive integer $K_2=K_2(k)$ such that, among any $K_2$ points to which $H$ has been assigned, we can find a subset $S$ of at least $k$ points for which: 
\begin{itemize}
    \item the three properties in the statement of Proposition~\ref{prop:first_structure} are satisfied;

    \item the sectors of $H$ with respect to every element of $S$ are all identical (including the order).
\end{itemize}
\end{proposition}

\begin{proof} By making $K_2$ sufficiently large, we may assume that there is a set $\tilde{S}$ of $10^{10}k$ points for which all three properties in the statement of Proposition~\ref{prop:first_structure} hold. Let $a$ and $b$ be the anchors of all assignments of $H$ to the elements of $\tilde S$, and label the elements of $\tilde S$ as $s_1,\dots,s_{10^{10}k}$ so that the third property in the statement of Proposition~\ref{prop:first_structure} holds. Without loss of generality, we assume that $\measuredangle as_ib<\pi$ for every $i$. 

For every $1\leq i\leq 10^{10}k$, each element $q\not\in H$ of the sector of $H$ with respect to $s_i$ will be classified as being either \textit{nearby} or \textit{away} with respect to $s_i$, depending on whether $q$ is contained in the convex hull of $H\cup\{p\}$ or not, respectively. By the first property in the statement of Proposition~\ref{prop:first_structure}, if some point $q\notin H$ is in the sector of $H$ with respect to two points $s_i$ and $s_j$, then it is either nearby with respect to both or away with respect to both. Since $\triangle s_iab$ is empty for every $s_i$, every point $q\notin H$ which is nearby with respect to $s_i$ can be seen to satisfy either $\measuredangle qs_ia<\pi$ or $\measuredangle bs_iq<\pi$ (but not both).

Suppose that the point $q\notin H$ is nearby with respect to some $s_i$ with $i>2$, and that $\measuredangle bs_iq<\pi$. Consider any $s_j$ with $j>i$. We claim that $q$ is nearby with respect to $s_j$, and that $\measuredangle bs_jq<\pi$. Notice that $q$ belongs to the convex hull of $H\cup\{s_i\}$. Thus, by the third property in the statement of Proposition~\ref{prop:first_structure}, it must be the case that $q\in\angle bs_is_j$. For the sake of contradiction, suppose that $\measuredangle qs_jb<\pi$ (or, equivalently, that $\measuredangle bs_jq>\pi$). Since $\triangle s_jab$ is empty, this can only occur if $q\in\angle s_is_ja$ (see Figure~\ref{fig:s1sisj}). As $q$ is in the sector of $H$ with respect to $s_i$, we can use the third property from Proposition~\ref{prop:first_structure} with $s_1$ and $s_i$ to conclude that $\measuredangle s_1s_jq<\pi$ and thus $s_1abqs_i$ is a convex pentagon (again, see Figure~\ref{fig:s1sisj}). Let $s_1'$ (resp. $q'$) denote the element of $P$ with the smallest non-zero Euclidean distance to $\overleftrightarrow{s_ia}$ (resp. $\overleftrightarrow{s_ib}$) that is contained in the triangle $\triangle s_1as_i$ (resp. $\triangle s_ibq$). Then, $s_1'abq's_i$ is a $5$-hole. This, however, contradicts our initial assumption that $H$ is assigned to $s_i$, because this $5$-hole would have been assigned to $s_i$ instead of $H$. Thus, $\measuredangle bs_jq<\pi$ must hold. It is not hard to see that $q$ must then belong to the sector of $H$ with respect to $s_j$, and thus must also be nearby with respect to this point. 

\begin{figure}[h]
    \centering
    
    \caption{The convex hull of $H\cup\{s_i\}$ has been shaded. The point $q$ must belong to the angular regions $\angle bs_is_j$ and $\angle s_is_ja$. Moreover, since $q$ is in the shaded region, the third property in Proposition~\ref{prop:first_structure} implies that $\measuredangle s_1s_is_j<\pi$. It is then easy to see that $s_1abqs_i$ (shown in red) is a convex pentagon.}
    \label{fig:s1sisj}
\end{figure}

By a symmetric argument, if $q\not\in H$ is nearby with respect to some $s_i$ with $i<10^{10}k$ and $\measuredangle qs_ia<\pi$, then $q$ is also nearby with respect to every $s_j$ with $j<i$, and it must be the case that $\measuredangle qs_ja<\pi$.

In particular, the above discussion implies that if $q\not\in H$ is nearby with respect to some $s_i$ with $1<i<10^{10}k$, then $q$ must be nearby with respect to $s_1$ or $s_{10^{10}k}$ (possibly both). Given that there are at most $5$ points which are nearby with respect to $s_1$ and at most $5$ points which are nearby with respect to $s_{10^{10}k}$, there is a set of at most $10$ points which contains all points $q$ that are nearby with respect to at least one element of $\tilde S$. For each $s_i$, given the sector $Q$ of $H$ with respect to $s_i$, we define the \textit{nearby sector} of $H$ with respect to $s_i$ as the ordered set obtained from $Q$ by ignoring all the away points. By the Pigeonhole-Principle, there is a subset $ S$ of $\tilde S$ of size at least $10^{10}k/(2^{10}10!)>k$ for which the corresponding nearby sectors of $H$ are all identical; let $p_1,\dots,p_t$ be the elements of these nearby sectors, in order. Note that $5\leq t\leq 10$, as the vertices of $H$ must lie among these points.

Next, we show that there is no index $1<i<t$ with $\measuredangle p_{i-1}p_ip_{i+1}>\pi$. Indeed, suppose that $\measuredangle p_{i-1}p_ip_{i+1}>\pi$ and let $s$ and $s'$ be two distinct elements of $S$. Then, without loss of generality, $sp_{i-1}p_ip_{i+1}s'$ is a convex pentagon, in that order. Furthermore, the convex quadrilaterals $sp_{i-1}p_ip_{i+1}$ and $s'p_{i-1}p_ip_{i+1}$ are both empty. Hence, if $s''\in P$ is the point in $\triangle sp_{i+1}s'$ with smallest non-zero Euclidean distance to $\overleftrightarrow{sp_{i+1}}$, then $sp_{i-1}p_ip_{i+1}s''$ is a $5$-hole which would have been assigned to $s$ instead of $H$ (see Figure~\ref{fig:pushing_from_top}a). This is a contradiction. Now, since $\measuredangle p_{i-1}p_ip_{i+1}>\pi$ for all $1<i<t$, the assignment of $H$ to $s_i$ must be either bad or dubious for every $s_i\in S$.

Finally, we claim that there are no away points in the sector of $H$ with respect to any of the elements of $S$, and thus all these sectors must be outright identical. Suppose, for the sake of contradiction, that there exists some point $q$ that is away with respect to some $s\in S$. Then, there exists some index $1\leq j\leq t-3$ with $q\in\angle p_jsp_{j+3}$, and it must be the case that $\measuredangle p_jqp_{j+3}>\pi$. Let $q'\in P$ be the point in $\triangle p_jqp_{j+3}$ with smallest non-zero Euclidean distance to $\overleftrightarrow{p_jp_{j+3}}$. Then, $p_jp_{j+1}p_{j+2}p_{j+3}q'$ is a $5$-hole (see Figure~\ref{fig:pushing_from_top}b). However, the assignment of this $5$-hole to $s$ would be good and trustworthy, while the assignment of $H$ to $s$ is dubious. This contradiction concludes the proof.
\end{proof} 

\end{document}
